\SourceLeaf{419}%
\noindent then
\[
(P\cos\varphi^2+2Q\cos\varphi\sin\varphi+R\sin\varphi^2)ds=(p\sin\varphi-q\cos\varphi)d\varphi,
\]
hence
\[
\frac1r=\frac{P\cos\varphi^2+2Q\cos\varphi\sin\varphi+R\sin\varphi^2}{p\sin\varphi-q\cos\varphi}
\]
or, since
\[
\cos\varphi=\frac{\mp q}{\BareRoot{(pp+qq)}},\qquad\sin\varphi=\frac{\pm p}{\BareRoot{(pp+qq)}}:
\]
\[
\pm\frac1r=\frac{Pqq-2Qpq+Rpp}{(pp+qq)^{\frac32}}.
\]

\Needspace{5\baselineskip}
\begin{center}4.\end{center}
\noindent\hspace*{1.2em}The ambiguous sign in the last formula may seem surprising at first, but on closer consideration it is entirely in order. For this expression depends only on the partial differentials of $V$, and the function $V$ in itself represents only the form of the curved line, without specifying the sense in which it is considered to be traced. It must therefore remain undecided whether the curve is convex towards the right or the left until this sense is determined from elsewhere. The same applies to determining $\varphi$ from its tangent, which by itself leaves an ambiguity of $180^{0}$. The sense in which the curve is traced can be specified in various ways.

I. By the sign of the change in $x$. If $x$ increases, $\cos\varphi$ must be positive. The upper signs therefore apply when $q$ has a negative value, and the lower when $q$ has a positive value. The reverse holds when $x$ decreases.

II. By the sign of the change in $y$. If $y$ increases, a positive value of $p$ requires the upper signs and a negative one the lower. The reverse holds when $y$ decreases.

III. By the sign of the value the function $V$ takes on when one leaves the curve. Let the variations of $x,y$ when one leaves the curve perpendicular to the tangent towards the right, hence in the direction $\varphi+90^{0}$, be $\delta x,\delta y$, with magnitude of the normal $=\delta\varrho$. We then\space%
